 This thesis focuses on the design and verification of an advanced control system for an unmanned underwater vehicle (UUV). While initially targeted at a remotely operated vehicle (ROV), the system features an open architecture that ensures seamless scalability toward fully autonomous underwater vehicle (AUV) applications. The core engineering challenge addressed in this research is the analytical derivation of a nonlinear, 6 degrees of freedom (6-DOF) mathematical model of the vehicle's dynamics, meticulously accounting for aquatic environment specificities such as added mass effects and viscous hydrodynamic damping. The study implements and compares classical control structures (PID) with custom control algorithms designed to compensate for the individual nonlinear dynamics of the drone. Due to the severe attenuation of electromagnetic waves underwater and the lack of GPS signals, the closed-loop control relies entirely on sensor fusion, utilizing data from an inertial measurement unit (IMU) and a piezoelectric pressure sensor. The control system software was developed within the ROS 2 framework and validated through Software-in-the-loop (SITL) simulations using the Gazebo physics engine. The entire software architecture has been engineered for direct deployment on physical hardware, specifically targeting a low-level Pixhawk flight controller running ArduSub, paired with a Raspberry Pi companion computer.